%% --- 6.1 Setup ---

%\subsection{Setup}
%\label{sec:exp:setup}

\paragraph{Benchmark.}
We evaluate on the Automated Design of Research Systems (ADRS) benchmark~\citep{cheng2025barbarians, cheng2025letthebarbarians}, which collects five research problems from computer systems: \textbf{Prism} (LLM-serving model placement across GPUs), \textbf{Cloudcast} (cloud network cost optimization), \textbf{EPLB} (expert-parallel load balancing for MoE models), \textbf{LLM-SQL} (tabular data layout for LLM prefix cache reuse), and \textbf{TXN} (transaction scheduling for makespan minimization).
Each task provides a fixed evaluator, starter code, and scoring metric.
We choose ADRS as our primary benchmark for three reasons: (1)~the tasks are drawn from real-world systems-optimization problems with established human baselines, (2)~the leaderboard provides both human expert and recent LLM-agent baselines, enabling apples-to-apples comparison, and (3)~the gold-standard evaluators are deterministic enough to support Score Verification and Specification Violation detection.
Previous studies observed that ADRS evaluators exhibit stochastic variance across runs~\citep{cemri2026adaevolve, liu2026evox}. We run each evaluator five times and compare against the adaptive tolerance $\max(1\%,\, 3\sigma/|\bar{s}|)$ to account for inherent evaluator variance.

\paragraph{Baseline Systems.}
We evaluate four baseline systems and \sys{}.
All baselines are open-source, enabling us to adapt each to the ADRS benchmark for controlled comparison under identical conditions.
They span the design spectrum from highly structured scaffolding to fully autonomous agents, providing coverage across the major paradigms for AI research agents.
\begin{itemize}[nosep,leftmargin=*]
  \item \textbf{Sakana AI-Scientist v2 (Sakana)}~\citep{yamada2025aiscientistv2}: Best-first tree search (BFTS) with a 4-stage experiment manager (preliminary investigation, hyperparameter tuning, research agenda execution, ablation studies) and a separate LLM writeup pipeline.
  \item \textbf{AutoResearchClaw (ARC)}~\citep{li2024autoresearchclaw}: 23-stage waterfall pipeline with multi-phase code generation (blueprint planning, sequential file generation, exec-fix loop, multi-agent review) and multi-source literature retrieval (OpenAlex, Semantic Scholar, arXiv, Google Scholar).
  \item \textbf{DeepScientist (DS)}~\citep{huang2025deepscientist}: Skill-based single-agent system on Codex CLI with separate code and write skills, using MCP tool servers for execution, memory, and artifacts.
  \item \textbf{AI-Researcher (AIR)}~\citep{tang2025airesearcher}: Orchestrated multi-agent system with specialized survey, coding, and writing agents. Experimentation uses a code-validate-refine loop.
  \item \textbf{\sys{}}: Full pipeline with evidence chain maintenance from problem framing through paper composition (\S\ref{sec:system}).
\end{itemize}

\paragraph{ADRS Adaptation.}
Each baseline was adapted to the ADRS benchmark with varying levels of source modification, from prompt-only changes (DS) to patching 16--19 source files (Sakana, AIR) to interface with ADRS task specifications, evaluators, and the NeurIPS~2026 paper template.
Sakana required the most extensive prompt-level rework: its default stage goals assume ML-training workflows (e.g., ``tune learning rates,'' ``introduce datasets from HuggingFace''), causing most initial runs to train neural networks instead of optimizing the target functions. A full rewrite of stage goals and 14 prompt locations was required before runs produced valid ADRS solutions (Appendix~\ref{app:baseline_adaptation}).
We standardized on Gemini~3.1 Pro as the backbone LLM across all systems for both solver code generation and paper writing.
To ensure best-effort performance, we configured generous iteration and timeout budgets: up to 20 solver iterations per task---6.7$\times$ the default for ARC, though most tasks converge well before this cap---and 2-hour code generation windows.
Runs that crashed due to infrastructure issues (API timeouts, rate limits, LaTeX compilation errors) were re-attempted with fresh state, up to 3~attempts per run. No run was re-attempted to improve solver scores.
For each system, we run 3~seeds per task, producing 15~papers per system and 75~papers total.
Of these, 16 required at least one retry due to infrastructure issues. All 75~runs ultimately produced solver code and a compiled paper, though artifact quality varies.
\coeaudit{} is applied identically to all systems.
Full adaptation details are provided in Appendix~\ref{app:baseline_adaptation}.
